% GENERATED FILE. Edit canonical lessons, then run npm run generate:books.
% canonical-source-hash: fnv1a64:32023b575ef287dd
% canonical-lessons: KA-C12-kutumba, KA-S121-letter-sa, KA-S143-digit-four

\chapter{Family}
\label{ch:ka-family}
\hlchaptermodality{12}

\begin{chapteropening}
\textbf{By the end of this chapter:} \emph{I can name my parents and my brothers and sisters in Kannada, choosing the right word for older or younger.}

Name \kn{ಅಪ್ಪ} and \kn{ಅಮ್ಮ}, then the four age-graded siblings \kn{ಅಣ್ಣ} · \kn{ತಮ್ಮ} · \kn{ಅಕ್ಕ} · \kn{ತಂಗಿ}, where age is chosen before gender.
\end{chapteropening}

\section[appa amma aṇṇa tamma akka taṅgi]{\kn{ಅಪ್ಪ}, \kn{ಅಮ್ಮ}, and four sibling words — the same system as Tamil}

\label{lesson:KA-C12-kutumba}



\begin{quote}
Kannada asks the same question Tamil does about any sibling: \textbf{older, or younger?} — and the words for answering it are almost identical.
\end{quote}

\begin{cousinweb}
\begin{itemize}
\raggedright
  \item \textbf{\kn{ಅಪ್ಪ}} (\emph{appa}) = \textquotedblleft{}\textbf{father}\textquotedblright{} — matches Tamil \emph{appā} closely.
  \item \textbf{\kn{ಅಮ್ಮ}} (\emph{amma}) = \textquotedblleft{}\textbf{mother}\textquotedblright{} — matches Tamil \emph{ammā} closely.
\end{itemize}
\end{cousinweb}

\begin{cousinweb}
\noindent
\begin{tabularx}{\linewidth}{@{}>{\raggedright\arraybackslash}X>{\raggedright\arraybackslash}X>{\raggedright\arraybackslash}X@{}}
\toprule
\textbf{Kannada} & \textbf{meaning} & \textbf{compare Tamil} \\
\midrule
\textbf{\kn{ಅಣ್ಣ}} (\emph{aṇṇa}) & \textbf{older} brother & \emph{aṇṇaṉ} — near match \\
\textbf{\kn{ತಮ್ಮ}} (\emph{tamma}) & \textbf{younger} brother & \emph{tambi} — near match \\
\textbf{\kn{ಅಕ್ಕ}} (\emph{akka}) & \textbf{older} sister & \emph{akkā} — near match \\
\textbf{\kn{ತಂಗಿ}} (\emph{tangi}) & \textbf{younger} sister & \emph{tangai} — near match \\
\bottomrule
\end{tabularx}

Just like Tamil, Kannada has \textbf{no single word} that simply means \textquotedblleft{}brother\textquotedblright{} or \textquotedblleft{}sister\textquotedblright{} — you name the sibling's \textbf{relative age} every time, not just their gender. This age-first system runs across the whole Dravidian family, not just one language.
\end{cousinweb}

\subsection*{Your turn}
\begin{itemize}
\raggedright
  \item \emph{Say it:} \textquotedblleft{}appa, amma\textquotedblright{} — father, mother
  \item \emph{Say it:} \textquotedblleft{}aṇṇa, tamma\textquotedblright{} — older/younger brother
  \item \emph{Say it:} \textquotedblleft{}akka, tangi\textquotedblright{} — older/younger sister
  \item \emph{Recall:} say \emph{dayaviṭṭu}, then read \textbf{\kn{ಕ್ಷಮಿಸಿ}}
\end{itemize}

\begin{tcolorbox}[breakable,colback=teal!4,colframe=teal!35!black,title={Before you move on}]
How many Kannada sibling words are there, and what do they split on? (\textbf{Four} — \kn{ಅಣ್ಣ}/\kn{ತಮ್ಮ}, \kn{ಅಕ್ಕ}/\kn{ತಂಗಿ} — split by \textbf{relative age}, not just gender.) How closely do Kannada's words match Tamil's? (\textbf{Very closely} — nearly the same forms across the whole set.)
\end{tcolorbox}
\section[sa]{\kn{ಸ} — one character, met inside words you already say}

\label{lesson:KA-S121-letter-sa}



\begin{quote}
Before the new one: \kn{ಹ} — what does it do?

One character this time. Just one — and you have been saying it for pages without knowing which mark on the page it was.
\end{quote}

\subsection*{Script you'll notice: \kn{ಸ}}
\textbf{\kn{ಸ}} — \emph{sa}.

It is a \textbf{consonant}, and in this script a consonant is never bare: it comes with an \emph{a} already in it. So it is not \emph{s}, it is \textbf{sa}.

You already say these, and every one of them has \kn{ಸ} somewhere inside it:

\begin{itemize}
\raggedright
  \item \textbf{\kn{ನಮಸ್ಕಾರ}} \emph{namaskāra} — hello / greetings (namaskāra — \textquotedblleft{}a making of a bow\textquotedblright{})
  \item \textbf{\kn{ಸರಿ}} \emph{sari} — okay / alright / correct (sari)
  \item \textbf{\kn{ಹೆಸರು}} \emph{hesaru} — name
  \item \textbf{\kn{ಸಂತೋಷ}} \emph{santōṣa} — joy / pleased to meet you
\end{itemize}

\subsection*{Writing: \kn{ಸ} — copy what you see}
Put your pen on \kn{ಸ} and follow its line. Copy the shape you can see — slowly, and larger than it is printed.

\begin{quote}
This book does not yet tell you \textbf{where to start the character or which way to travel}. That is a real thing, taught with real variation from school to school, and it is not written down here until it can be written down with a source. Copying what is in front of you needs no such source, and it is how the shape gets into your hand in the meantime.
\end{quote}

\subsection*{Your turn}
\begin{itemize}
\raggedright
  \item \emph{Look:} at these words, and find \kn{ಸ} in the ones that have it
\end{itemize}

\begin{quote}
\kn{ನಮಸ್ಕಾರ}  ·  \kn{ಸರಿ}  ·  \kn{ಧನ್ಯವಾದ}
\end{quote}

\begin{itemize}
\raggedright
  \item \emph{Trace it:} \kn{ಸ} three times, saying \emph{sa} as you finish each one
  \item \emph{Look:} back at any page of this chapter and find \kn{ಸ} once more
\end{itemize}

\begin{tcolorbox}[breakable,colback=teal!4,colframe=teal!35!black,title={Before you move on}]
Which character is this — \kn{ಸ}? What sound does it carry? (\emph{\textbf{sa}}.) Name one word you already say that contains it.
\end{tcolorbox}
\section[4]{\kn{೪} — the digit for nālku}

\label{lesson:KA-S143-digit-four}



\begin{quote}
Before the new shape: \textbf{\kn{೩}} — which number is it?

One digit this time, and you already know its name.
\end{quote}

\subsection*{Script you'll notice: \kn{೪}}
\textbf{\kn{೪}} — \textbf{4}, the digit for \textbf{\kn{ನಾಲ್ಕು}} (\emph{nālku}), \textquotedblleft{}four\textquotedblright{}.

You have now met four of the ten. The words came first, when you learned to count; the shapes are arriving one at a time, on top of words you can already say.

\subsection*{Writing: \kn{೪} — copy what you see}
Put your pen on \kn{೪} and follow its line. Copy the shape in front of you — slowly, and larger than it is printed.

\begin{quote}
This book does not yet tell you \textbf{where to start the character or which way to travel}. That is taught with real variation from school to school, and it is not written down here until it can be written down with a source. Copying what you can see needs no such source.
\end{quote}

\subsection*{Your turn}
\begin{itemize}
\raggedright
  \item \emph{Look:} at the digits you have met, and name each one aloud
\end{itemize}

\begin{quote}
\kn{೧}  ·  \kn{೨}  ·  \kn{೩}  ·  \kn{೪}
\end{quote}

\begin{itemize}
\raggedright
  \item \emph{Say it:} \textquotedblleft{}nālku\textquotedblright{} as you point at \kn{೪}
\end{itemize}

\begin{tcolorbox}[breakable,colback=teal!4,colframe=teal!35!black,title={Before you move on}]
Which number is \textbf{\kn{೪}}? (\textbf{4} — \emph{nālku}, \kn{ನಾಲ್ಕು}.) Name every digit you have met so far, in order.
\end{tcolorbox}
